\appendix

\section{Questionnaire : Usages de l'\gls{ia} générative à Polytech}
\label{annexe:questionnaire}

\subsection*{Note méthodologique}

Ce questionnaire vise à documenter les usages réels de l'\gls{ia} générative au sein de Polytech Annecy-Chambéry, afin d'adapter les recommandations institutionnelles aux besoins spécifiques de chaque public (étudiants, enseignants-chercheurs, personnels administratifs et techniques). Quatre objectifs structurent cette enquête : (1) comparer les pratiques locales aux données nationales (n=30~000 répondants de l'enquête ministérielle), (2) identifier les besoins d'accompagnement différenciés par public, (3) repérer les champions \gls{ia} potentiels au sein des départements, et (4) documenter les enjeux spécifiques aux écoles d'ingénieurs (confidentialité des données industrielles, usage du code assisté par \gls{ia}, spécificités disciplinaires).

La méthodologie repose sur trois piliers validés scientifiquement. Les items du questionnaire national sont repris intégralement pour garantir la comparabilité des résultats (sections A à E, tronc commun). Les échelles psychométriques validées internationalement sont intégrées : échelle d'usage académique de ChatGPT (Abbas et al., 2024, 8 items, $\alpha>0.70$), modèle UTAUT2 adapté à l'\gls{ia} (5 construits), et échelle d'auto-efficacité GSE-6AI (6 items). Enfin, des questions spécifiques aux écoles d'ingénieurs complètent le dispositif (section F, différenciée par public)~\cite{pascal2025ia,cge2024enquete,abbas2024harmful}.

Le questionnaire adopte une architecture modulaire : un tronc commun de 42 items (sections A à E) pour tous les répondants, suivi d'une branche spécifique de 8 items selon le profil déclaré (section F), et d'une clôture commune (section G). La durée estimée de réponse est de 10 à 13 minutes. L'anonymat complet est garanti : aucune donnée personnelle identifiante n'est collectée. Les données agrégées serviront exclusivement à l'élaboration du plan d'action institutionnel. Un engagement de restitution des résultats aux participants est pris.

\bigskip

\subsection{Section A — Profil et contexte}

\textbf{A1.} Vous êtes :
\begin{itemize}[label=$\square$, leftmargin=1.5cm]
  \item Étudiant·e
  \item Enseignant·e-chercheur·euse
  \item Personnel administratif, technique, bibliothèque, ingénierie, social et santé (BIATSS)
\end{itemize}

\textbf{A2.} \textit{[Si Étudiant]} Dans quelle filière êtes-vous inscrit·e ?
\begin{itemize}[label=$\square$, leftmargin=1.5cm]
  \item Écologie Industrielle et Territoriale (EIT)
  \item INformatique Données Usages (IDU)
  \item Mécanique Mécatronique et Matériaux (MM)
  \item Système Embarqué et Automatisme (SEA)
  \item Cycle préparatoire (PeiP)
  \item Autre : \underline{\hspace{4cm}}
\end{itemize}

\textbf{A3.} \textit{[Si Étudiant]} Vous êtes en :
\begin{itemize}[label=$\square$, leftmargin=1.5cm]
  \item Formation Initiale
  \item Alternance
  \item Formation Continue
  \item Autre : \underline{\hspace{4cm}}
\end{itemize}

\textbf{A4.} \textit{[Si Étudiant]} Vous êtes en :
\begin{itemize}[label=$\square$, leftmargin=1.5cm]
  \item PeiP (1\textsuperscript{re} ou 2\textsuperscript{e} année cycle préparatoire)
  \item 3\textsuperscript{e} année (1\textsuperscript{re} année cycle ingénieur)
  \item 4\textsuperscript{e} année (2\textsuperscript{e} année cycle ingénieur)
  \item 5\textsuperscript{e} année (3\textsuperscript{e} année cycle ingénieur)
\end{itemize}


\textbf{A5.} \textit{[Si Enseignant-chercheur]} Quelle est votre discipline principale d'enseignement ?
\begin{itemize}[label=$\square$, leftmargin=1.5cm]
  \item Informatique et sciences du numérique
  \item Mathématiques et modélisation
  \item Physique et sciences de l'ingénieur
  \item Mécanique et matériaux
  \item Environnement et énergétique
  \item Sciences humaines et sociales, langues
  \item Autre : \underline{\hspace{4cm}}
\end{itemize}

\textbf{A6.} \textit{[Si Enseignant-chercheur]} Vous êtes rattaché·e à :
\begin{itemize}[label=$\square$, leftmargin=1.5cm]
  \item LISTIC
  \item SYMME
  \item LOCIE
  \item IREGE
  \item LAPP
  \item LAMA
  \item Autre : \underline{\hspace{4cm}}
\end{itemize}

\textbf{A7.} \textit{[Si BIATSS]} Quelle est votre fonction principale ?
\begin{itemize}[label=$\square$, leftmargin=1.5cm]
  \item Administration (gestion, finances, RH)
  \item Scolarité et relations étudiantes
  \item Informatique et systèmes d'information
  \item Ingénierie pédagogique
  \item Bibliothèque et documentation
  \item Technique et logistique
  \item Autre : \underline{\hspace{4cm}}
\end{itemize}

\textbf{A8.} Saviez-vous qu'une charte sur l'usage de l'\gls{ia} générative a récemment été adoptée par l'\gls{usmb} ?
\begin{itemize}[label=$\square$, leftmargin=1.5cm]
  \item Oui, j'en avais déjà connaissance
  \item J'en avais vaguement entendu parler, sans en connaître le contenu
  \item Non, je l'apprends à l'instant
\end{itemize}

\textbf{A9.} Estimez-vous que cette charte vous donne
des orientations suffisamment claires pour guider votre
usage de l'\gls{ia} générative ?
\begin{itemize}[label=$\square$, leftmargin=1.5cm]
  \item Oui, elle est claire et suffisante
  \item Partiellement, certains points restent flous
  \item Non, elle est trop vague ou incomplète
  \item Je ne suis pas en mesure de me prononcer
\end{itemize}

\medskip

\subsection{Section B — Usage actuel de l'\gls{ia} générative}

\textbf{B1.} Utilisez-vous des outils d'\gls{ia} générative (ChatGPT, Claude, Gemini, Copilot, etc.) ?
\begin{itemize}[label=$\square$, leftmargin=1.5cm]
  \item Oui
  \item Non $\rightarrow$ \textit{Passez à la question B9}
\end{itemize}

\textbf{B2.} Depuis combien de temps utilisez-vous l'\gls{ia} générative ?
\begin{itemize}[label=$\square$, leftmargin=1.5cm]
  \item Moins de 3 mois
  \item Entre 3 et 6 mois
  \item Entre 6 mois et 1 an
  \item Entre 1 an et 2 ans
  \item Plus de 2 ans
\end{itemize}

\textbf{B3.} À quelle fréquence utilisez-vous l'\gls{ia} générative ?
\begin{itemize}[label=$\square$, leftmargin=1.5cm]
  \item Jamais (je ne l'utilise plus)
  \item Rarement (moins d'une fois par mois)
  \item Occasionnellement (1 à 3 fois par mois)
  \item Régulièrement (1 à 2 fois par semaine)
  \item Fréquemment (plusieurs fois par semaine)
  \item Quotidiennement
\end{itemize}

\textbf{B4.} Dans quel(s) contexte(s) utilisez-vous l'\gls{ia} générative ? \textit{(Plusieurs réponses possibles)}
\begin{itemize}[label=$\square$, leftmargin=1.5cm]
  \item \textit{[Étudiant]} Pour mes études (travaux, projets, révisions)
  \item \textit{[Enseignant]} Pour préparer mes cours et mes supports pédagogiques
  \item \textit{[Enseignant]} Pour mes activités de recherche
  \item \textit{[BIATSS]} Pour mes tâches professionnelles administratives
  \item Usage personnel (hors cadre professionnel/études)
\end{itemize}

\textbf{B5.} Quel est votre niveau de compétence auto-évalué dans l'utilisation de l'\gls{ia} générative ?
\begin{itemize}[label=$\square$, leftmargin=1.5cm]
  \item Débutant : Usage basique (Je pose des questions simples, je l'utilise un peu comme un moteur de recherche).
  \item Intermédiaire : Usage actif (Je sais rédiger des consignes détaillées pour générer ou corriger du texte, du code, ou résumer des documents).
  \item Avancé : Usage automatisé (J'intègre l'IA dans mes méthodes de travail pour automatiser des tâches répétitives ou analyser de gros volumes de données).
  \item Expert : Usage technique ou de référence (Je développe des outils complexes avec l'IA, ou je suis la personne ressource qui aide et forme les autres).
\end{itemize}

\textbf{B6.} Avez-vous déjà déclaré auprès d'un enseignant votre usage de l'\gls{ia} générative pour un travail académique ?
\begin{itemize}[label=$\square$, leftmargin=1.5cm]
  \item Oui, toujours
  \item Oui, parfois
  \item Non, jamais
  \item Non concerné (je ne l'utilise pas pour les travaux académiques)
\end{itemize}

\textbf{B7.} Avez-vous reçu une formation à l'usage de l'\gls{ia} générative ?
\begin{itemize}[label=$\square$, leftmargin=1.5cm]
  \item Oui, une formation institutionnelle (Polytech, \gls{usmb})
  \item Oui, une formation externe
  \item Non, je suis autodidacte
  \item Non, je n'ai reçu aucune formation
\end{itemize}

\textbf{B8.} Seriez-vous intéressé·e par une formation sur l'usage éthique et efficace de l'\gls{ia} générative ?
\begin{itemize}[label=$\square$, leftmargin=1.5cm]
  \item Oui, très intéressé·e
  \item Oui, plutôt intéressé·e
  \item Non, peu intéressé·e
  \item Non, pas du tout intéressé·e
\end{itemize}

\textbf{B9.} \textit{[Si non-utilisateur]} Quelle est la principale raison pour laquelle vous n'utilisez pas l'\gls{ia} générative ? \textit{(Une seule réponse)}
\begin{itemize}[label=$\square$, leftmargin=1.5cm]
  \item Je ne connais pas ces outils
  \item Je ne vois pas d'utilité pour mon travail/mes études
  \item Je ne sais pas comment les utiliser
  \item Je crains les problèmes d'intégrité académique ou éthique
  \item Je ne fais pas confiance à la fiabilité des résultats
  \item L'établissement n'a pas donné de règles claires
  \item Je n'ai pas accès à ces outils adaptés (outils payants, etc)
  \item Autre raison : \underline{\hspace{4cm}}
\end{itemize}

\medskip

\subsection{Section C — Outils et tâches}

\textbf{C1.} Quels outils d'\gls{ia} générative utilisez-vous ? \textit{(Plusieurs réponses possibles)}
\begin{itemize}[label=$\square$, leftmargin=1.5cm]
  \item ChatGPT (OpenAI)
  \item Claude (Anthropic)
  \item Google Gemini
  \item Microsoft Copilot (Bing Chat)
  \item le Chat (Mistral AI)
  \item GitHub Copilot, Cursor, Codex, Claude Code (Génération de code)
  \item Perplexity (Recherche de documentaires)
  \item Modèles locaux (Ollama, vLLM, etc)
  \item Autre : \underline{\hspace{4cm}}
\end{itemize}

\textbf{C2.} Pour quelles tâches utilisez-vous l'\gls{ia} générative ? \textit{(Plusieurs réponses possibles)}
\begin{itemize}[label=$\square$, leftmargin=1.5cm]
  \item Recherche d'information et documentation
  \item Rédaction de textes (rapports, courriels, synthèses)
  \item Résumé de documents longs
  \item Traduction de textes
  \item Génération ou complétion de code informatique
  \item Débogage de code
  \item Clarification de concepts théoriques
  \item Génération d'exercices ou de questions d'examen
  \item Création de supports visuels (diagrammes, schémas)
  \item Aide à la correction ou à l'évaluation de travaux
  \item Brainstorming et génération d'idées
  \item Analyse de données (tableur, CSV, etc.)
  \item Autre : \underline{\hspace{4cm}}
\end{itemize}

\textbf{C3.} Quelle proportion de votre production finale est générée directement par l'\gls{ia} ?
\begin{itemize}[label=$\square$, leftmargin=1.5cm]
  \item 0\% (je n'utilise l'\gls{ia} que pour des idées, pas pour la rédaction finale)
  \item Moins de 25\%
  \item Entre 25\% et 50\%
  \item Entre 50\% et 75\%
  \item Plus de 75\%
\end{itemize}

\textbf{C4.} Vérifiez-vous systématiquement les informations fournies par l'\gls{ia} générative ?
\begin{itemize}[label=$\square$, leftmargin=1.5cm]
  \item Toujours
  \item Souvent
  \item Parfois
  \item Rarement
  \item Jamais
\end{itemize}

\medskip

\subsection{Section D — Perceptions et attitudes}

Les affirmations suivantes utilisent une échelle de Likert en 5 points : 1 = Pas du tout d'accord ; 2 = Plutôt pas d'accord ; 3 = Neutre ; 4 = Plutôt d'accord ; 5 = Tout à fait d'accord.

\textbf{D1.} L'utilisation de l'\gls{ia} générative améliore ma productivité.
\begin{itemize}[label=$\square$, leftmargin=1.5cm]
  \item 1 \quad $\square$ 2 \quad $\square$ 3 \quad $\square$ 4 \quad $\square$ 5
\end{itemize}

\textbf{D2.} Apprendre à utiliser l'\gls{ia} générative est facile pour moi.
\begin{itemize}[label=$\square$, leftmargin=1.5cm]
  \item 1 \quad $\square$ 2 \quad $\square$ 3 \quad $\square$ 4 \quad $\square$ 5
\end{itemize}

\textbf{D3.} Je fais confiance aux informations fournies par l'\gls{ia} générative.
\begin{itemize}[label=$\square$, leftmargin=1.5cm]
  \item 1 \quad $\square$ 2 \quad $\square$ 3 \quad $\square$ 4 \quad $\square$ 5
\end{itemize}

\textbf{D4.} Mes enseignants/collègues encouragent l'utilisation de l'\gls{ia} générative.
\begin{itemize}[label=$\square$, leftmargin=1.5cm]
  \item 1 \quad $\square$ 2 \quad $\square$ 3 \quad $\square$ 4 \quad $\square$ 5
\end{itemize}

\textbf{D5.} J'ai l'intention d'utiliser davantage l'\gls{ia} générative à l'avenir.
\begin{itemize}[label=$\square$, leftmargin=1.5cm]
  \item 1 \quad $\square$ 2 \quad $\square$ 3 \quad $\square$ 4 \quad $\square$ 5
\end{itemize}

\textbf{D6.} Je suis confiant·e dans ma capacité à utiliser l'\gls{ia} générative efficacement.
\begin{itemize}[label=$\square$, leftmargin=1.5cm]
  \item 1 \quad $\square$ 2 \quad $\square$ 3 \quad $\square$ 4 \quad $\square$ 5
\end{itemize}

\textbf{D7.} Je sais comment vérifier la fiabilité des réponses de l'\gls{ia} générative.
\begin{itemize}[label=$\square$, leftmargin=1.5cm]
  \item 1 \quad $\square$ 2 \quad $\square$ 3 \quad $\square$ 4 \quad $\square$ 5
\end{itemize}

\textbf{D8.} L'\gls{ia} générative représente une opportunité pour mon travail/mes études.
\begin{itemize}[label=$\square$, leftmargin=1.5cm]
  \item 1 \quad $\square$ 2 \quad $\square$ 3 \quad $\square$ 4 \quad $\square$ 5
\end{itemize}

\textbf{D9.} L'\gls{ia} générative représente une menace pour mon travail/mes études.
\begin{itemize}[label=$\square$, leftmargin=1.5cm]
  \item 1 \quad $\square$ 2 \quad $\square$ 3 \quad $\square$ 4 \quad $\square$ 5
\end{itemize}

\textbf{D10.} Je suis préoccupé·e par les inexactitudes (hallucinations) de l'\gls{ia} générative.
\begin{itemize}[label=$\square$, leftmargin=1.5cm]
  \item 1 \quad $\square$ 2 \quad $\square$ 3 \quad $\square$ 4 \quad $\square$ 5
\end{itemize}

\medskip

\subsection{Section E — Freins et besoins}

\textbf{E1.} Quels sont les principaux freins à votre utilisation de l'\gls{ia} générative ? \textit{(Plusieurs réponses possibles)}
\begin{itemize}[label=$\square$, leftmargin=1.5cm]
  \item Manque de formation
  \item Incertitude sur les règles d'usage autorisées
  \item Crainte de problèmes d'intégrité académique ou déontologique
  \item Manque de confiance dans la fiabilité des résultats
  \item Préoccupations sur la confidentialité des données
  \item Temps nécessaire pour apprendre à utiliser ces outils
  \item Absence de soutien de ma hiérarchie ou de mes enseignants
  \item Coût des versions payantes
  \item Impact environnemental
  \item Autre : \underline{\hspace{4cm}}
\end{itemize}

\textbf{E2.} Quels types de formation souhaiteriez-vous recevoir sur l'\gls{ia} générative ? \textit{(Plusieurs réponses possibles)}
\begin{itemize}[label=$\square$, leftmargin=1.5cm]
  \item Introduction générale aux outils disponibles
  \item Ingénierie de prompts (formulation de requêtes efficaces)
  \item Usages éthiques et respect de l'intégrité académique
  \item Vérification de la fiabilité des résultats
  \item Applications spécifiques à ma discipline/mon métier
  \item Sécurité et confidentialité des données
  \item Impact environnemental et sobriété numérique
  \item Autre : \underline{\hspace{4cm}}
\end{itemize}

\textbf{E3.} Quelle modalité de formation préféreriez-vous ?
\begin{itemize}[label=$\square$, leftmargin=1.5cm]
  \item Formation en ligne asynchrone (MOOC, vidéos)
  \item Ateliers en présentiel (2-4 heures)
  \item Formation par les pairs (collègues/étudiants référents)
  \item Ressources documentaires (guides, fiches pratiques)
  \item Webinaires interactifs
  \item Autre : \underline{\hspace{4cm}}
\end{itemize}

\textbf{E4.} Combien de temps seriez-vous prêt·e à consacrer chaque année à une formation sur l'\gls{ia} ?
\begin{itemize}[label=$\square$, leftmargin=1.5cm]
  \item Moins de 2 heures
  \item 2 à 5 heures
  \item 5 à 10 heures
  \item Plus de 10 heures
\end{itemize}
 
\medskip

\subsection{Section F — Questions spécifiques}

\subsubsection{Version Étudiants}



\textbf{F1.} Utilisez-vous l'\gls{ia} générative pour les tâches
d'apprentissage suivantes ? \textit{(Plusieurs réponses possibles)}
\begin{itemize}[label=$\square$, leftmargin=1.5cm]
  \item Comprendre des concepts ou des formules complexes du cours
  \item Synthétiser les points clés et préparer des fiches de révision
  \item Générer des questions d'entraînement et corriger mes réponses
  \item Rédiger des rapports, devoirs ou dissertations
  \item Lire et comprendre des articles ou documents techniques en langue
        étrangère
  \item Préparer des soutenances orales ou des présentations
  \item Je n'utilise pas l'\gls{ia} pour ces tâches
\end{itemize}


\textbf{F2.} Estimez-vous que l'émergence de l'\gls{ia} a modifié votre
perception des compétences fondamentales de votre spécialité ?
\begin{itemize}[label=$\square$, leftmargin=1.5cm]
  \item Oui, je peux me concentrer davantage sur la compréhension
        disciplinaire plutôt que sur la mémorisation de détails ou de syntaxe
  \item Oui, mais je crains de devenir trop dépendant·e de l'\gls{ia} au
        détriment de mes compétences de base
  \item Non, je ne perçois pas de changement notable
  \item Je n'ai pas encore réfléchi à cette question
\end{itemize}


\textbf{F3.} Vos cours ou projets impliquent-ils de la programmation ou
de l'écriture de code ?
\begin{itemize}[label=$\square$, leftmargin=1.5cm]
  \item Oui \hspace{1cm} 
  \item Non \hspace{0.9cm} $\rightarrow$ \textit{Passez directement à F9}
\end{itemize}

\textbf{F4.} Quels outils d'\gls{ia} utilisez-vous pour la programmation ?
\textit{(Plusieurs réponses possibles)}
\begin{itemize}[label=$\square$, leftmargin=1.5cm]
  \item ChatGPT / Claude / Gemini (mode conversationnel)
  \item GitHub Copilot (autocomplétion dans l'éditeur)
  \item Cursor / Windsurf (éditeur intégrant l'\gls{ia})
  \item Claude Code / Devin ou autre outil en mode agent (CLI)
  \item Autre : \underline{\hspace{4cm}}
\end{itemize}
 
\textbf{F5.} Lorsque vous utilisez un assistant de programmation \gls{ia},
quel mode d'interaction utilisez-vous principalement ?
\begin{itemize}[label=$\square$, leftmargin=1.5cm]
  \item Mode conversationnel : je pose des questions dans une interface de
        chat et je copie-colle le code manuellement
  \item Mode éditeur intégré : l'\gls{ia} propose des suggestions ou
        complétions directement dans mon éditeur de code
  \item Mode agent : je décris une tâche et l'\gls{ia} l'exécute de façon
        autonome (ex. Claude Code CLI)
  \item Mode mixte : j'utilise plusieurs des modes ci-dessus
\end{itemize}
 
\textbf{F6.} Dans vos travaux de programmation, quelle proportion du code
est générée par l'\gls{ia} ?
\begin{itemize}[label=$\square$, leftmargin=1.5cm]
  \item Quasi nulle ($< 10\,\%$)
  \item Minoritaire ($10$--$30\,\%$)
  \item Significative ($30$--$60\,\%$)
  \item Majoritaire ($60$--$90\,\%$)
  \item Quasi totale ($> 90\,\%$)
\end{itemize}
 
\textbf{F7.} Comment décrivez-vous généralement une tâche de programmation
à l'\gls{ia} ? \textit{(Plusieurs réponses possibles)}
\begin{itemize}[label=$\square$, leftmargin=1.5cm]
  \item Une phrase ou une description courte et simple
  \item Du code existant accompagné d'une demande de modification précise
  \item Un document de spécifications détaillé (Markdown, cahier des
        charges structuré)
  \item Une description de l'architecture globale pour que l'\gls{ia}
        génère le squelette du projet
  \item Des fichiers ou captures d'écran existants pour contextualiser le
        projet
\end{itemize}
 
\textbf{F8.} Une fois le code généré par l'\gls{ia}, que faites-vous ?
\begin{itemize}[label=$\square$, leftmargin=1.5cm]
  \item Je l'utilise directement sans vérification
  \item Je le parcours rapidement avant de l'utiliser
  \item Je le lis attentivement et comprends chaque ligne avant de
        l'utiliser
  \item Je le valide à l'aide de tests avant de l'utiliser
  \item Je le comprends, le modifie et n'utilise que ma version remaniée
\end{itemize}

\textbf{F9.} Avez-vous déjà uploadé des données issues de projets
industriels ou de stages dans un outil d'\gls{ia} générative public
(ChatGPT, Claude, etc.) ?
\begin{itemize}[label=$\square$, leftmargin=1.5cm]
  \item Oui
  \item Non
  \item Je ne sais pas si cela représente un risque
  \item Non concerné·e (je n'ai pas encore de projet industriel ou de
        stage)
\end{itemize}
 
\textbf{F10.} Êtes-vous conscient·e des risques liés à la confidentialité
lorsque vous utilisez des outils d'\gls{ia} publics avec des données de
projets industriels ?
\begin{itemize}[label=$\square$, leftmargin=1.5cm]
  \item Oui, tout à fait
  \item Oui, plutôt
  \item Non, plutôt pas
  \item Non, pas du tout
\end{itemize}

\textbf{F11.} Durant votre stage ou alternance, votre entreprise a-t-elle
défini des règles ou une politique claire concernant l'utilisation de
l'\gls{ia} générative ?
\begin{itemize}[label=$\square$, leftmargin=1.5cm]
  \item Oui, une charte ou des directives strictes encadrent son utilisation
  \item Oui, mais les règles sont floues ou peu appliquées en pratique
  \item Non, il n'y a aucune règle établie à ma connaissance
  \item Non concerné·e
\end{itemize}

\textbf{F12.} Votre entreprise met-elle à disposition des abonnements ou
des versions professionnelles sécurisées d'outils d'\gls{ia} (ex. GitHub
Copilot Enterprise, ChatGPT Team/Enterprise) ?
\begin{itemize}[label=$\square$, leftmargin=1.5cm]
  \item Oui, l'entreprise fournit et finance ces outils officiellement
  \item Non, je dois utiliser mes propres outils (comptes personnels gratuits ou payants)
  \item Non, et l'utilisation d'outils d'\gls{ia} tiers est techniquement bloquée ou formellement interdite
  \item Non concerné·e
\end{itemize}

\textbf{F13.} Comment percevez-vous l'attitude générale de votre encadrement
ou de votre entreprise face à l'\gls{ia} générative ?
\begin{itemize}[label=$\square$, leftmargin=1.5cm]
  \item Très favorable : son usage est activement encouragé pour gagner en productivité
  \item Tolérante : son usage est accepté, mais pas particulièrement mis en avant
  \item Méfiante : des inquiétudes (sécurité, qualité du code, éthique) freinent son adoption
  \item Non concerné·e
\end{itemize}
 
\textbf{F14.} Considérez-vous que l'usage de l'\gls{ia} générative pour
rédiger un devoir est une forme de triche ?
\begin{itemize}[label=$\square$, leftmargin=1.5cm]
  \item Oui, dans tous les cas
  \item Oui, si l'enseignant ne l'a pas explicitement autorisé
  \item Non, à condition de déclarer son usage
  \item Non, jamais
  \item Je ne sais pas
\end{itemize}
 
\textbf{F15.} Avez-vous déjà ressenti de l'anxiété ou de la culpabilité
à utiliser l'\gls{ia} générative pour vos études ?
\begin{itemize}[label=$\square$, leftmargin=1.5cm]
  \item Oui, souvent
  \item Oui, parfois
  \item Non, rarement
  \item Non, jamais
\end{itemize}

\subsubsection{Version Enseignants-chercheurs}

\textbf{F1.} Utilisez-vous l'\gls{ia} générative pour préparer vos enseignements ?
\begin{itemize}[label=$\square$, leftmargin=1.5cm]
  \item Oui, régulièrement
  \item Oui, occasionnellement
  \item Non, jamais
\end{itemize}

\textbf{F2.} \textit{[Si oui]} Pour quelles tâches pédagogiques ? \textit{(Plusieurs réponses possibles)}
\begin{itemize}[label=$\square$, leftmargin=1.5cm]
  \item Préparation de supports de cours
  \item Création d'exercices ou de questions d'examen
  \item Génération d'exemples ou de cas d'étude
  \item Aide à la correction de travaux étudiants
  \item Traduction de ressources pédagogiques
  \item Autre : \underline{\hspace{4cm}}
\end{itemize}

\textbf{F3.} Utilisez-vous l'\gls{ia} générative dans vos activités de recherche ?
\begin{itemize}[label=$\square$, leftmargin=1.5cm]
  \item Oui, régulièrement
  \item Oui, occasionnellement
  \item Non, jamais
\end{itemize}

\textbf{F4.} \textit{[Si oui]} Pour quelles activités de recherche ? \textit{(Plusieurs réponses possibles)}
\begin{itemize}[label=$\square$, leftmargin=1.5cm]
  \item Revue de littérature et veille scientifique
  \item Rédaction d'articles ou de rapports
  \item Génération ou débogage de code de traitement de données
  \item Traduction d'articles
  \item Aide à la formulation d'hypothèses
  \item Autre : \underline{\hspace{4cm}}
\end{itemize}

\textbf{F5.} Estimez-vous que vos étudiants utilisent l'\gls{ia} générative pour leurs travaux ?
\begin{itemize}[label=$\square$, leftmargin=1.5cm]
  \item Oui, la majorité (plus de 70\%)
  \item Oui, environ la moitié (40-70\%)
  \item Oui, une minorité (10-40\%)
  \item Non, très peu (moins de 10\%)
  \item Je ne sais pas
\end{itemize}

\textbf{F6.} Quelle est votre position concernant l'usage de l'\gls{ia}
générative par vos étudiants dans le cadre de vos cours ?
\begin{itemize}[label=$\square$, leftmargin=1.5cm]
  \item Interdit, je l'indique explicitement dans les règles du cours
  \item Autorisé sous conditions (déclaration obligatoire ou autorisation
        préalable)
  \item Libre, à condition de mentionner l'usage de l'\gls{ia}
  \item Totalement libre, sans restriction particulière
  \item Je n'ai pas encore de position ou de règle clairement définie
\end{itemize}

\textbf{F7.} Avez-vous communiqué explicitement vos règles d'usage de
l'\gls{ia} à vos étudiants en début de cours ?
\begin{itemize}[label=$\square$, leftmargin=1.5cm]
  \item Oui, par écrit (syllabus, règlement, consignes formalisées)
  \item Oui, oralement lors d'une séance
  \item Non, mais j'envisage de le faire
  \item Non, je ne pense pas que ce soit nécessaire
\end{itemize}

\textbf{F8.} Avez-vous adapté vos modalités d'évaluation en raison de l'\gls{ia} générative ?
\begin{itemize}[label=$\square$, leftmargin=1.5cm]
  \item Oui, substantiellement
  \item Oui, légèrement
  \item Non, pas encore
  \item Non, je ne pense pas que ce soit nécessaire
\end{itemize}

\textbf{F9.} Seriez-vous prêt·e à devenir référent·e \gls{ia} pour accompagner vos collègues ?
\begin{itemize}[label=$\square$, leftmargin=1.5cm]
  \item Oui, volontiers
  \item Oui, si je reçois une formation préalable
  \item Non, je ne me sens pas compétent·e
  \item Non, je n'ai pas le temps
\end{itemize}

\textbf{F10.} Utilisez-vous des outils de détection de contenu généré par
l'\gls{ia} (Turnitin, GPTZero, etc.) pour évaluer les travaux de vos
étudiants ?
\begin{itemize}[label=$\square$, leftmargin=1.5cm]
  \item Oui, de manière systématique
  \item Oui, occasionnellement
  \item Non, mais j'y ai réfléchi
  \item Non, je considère ces outils peu fiables
  \item Non, je n'estime pas cela nécessaire
\end{itemize}
 
\textbf{F11.} Selon vous, quel sera l'impact de l'\gls{ia} générative sur
le métier d'enseignant-chercheur ?
\begin{itemize}[label=$\square$, leftmargin=1.5cm]
  \item Une menace significative : certaines fonctions pédagogiques
        pourraient être remplacées
  \item Une transformation profonde, mais le rôle de l'enseignant évoluera
        sans disparaître
  \item Un impact limité : les compétences fondamentales de l'enseignement
        ne peuvent être remplacées
  \item Une opportunité : l'\gls{ia} libérera du temps pour des tâches à
        plus haute valeur ajoutée
  \item Je n'ai pas encore de position tranchée sur ce sujet
\end{itemize}
 
\textbf{F12.} Abordez-vous activement avec vos étudiants la question d'un
usage responsable de l'\gls{ia} générative ?
\begin{itemize}[label=$\square$, leftmargin=1.5cm]
  \item Oui, c'est intégré dans le contenu de mon cours
  \item Oui, j'en parle ponctuellement
  \item Non, je considère que c'est du ressort de l'établissement
  \item Non, cela ne relève pas de mes responsabilités pédagogiques
\end{itemize}

\subsubsection{Version BIATSS}

\textbf{F1.} Quelles sont les tâches qui composent principalement votre
travail au quotidien ? \textit{(Plusieurs réponses possibles)}
\begin{itemize}[label=$\square$, leftmargin=1.5cm]
  \item Rédaction ou traitement de courriels, courriers et communications
  \item Rédaction de rapports, comptes rendus ou documents administratifs
  \item Utilisation de tableurs (Excel) pour le traitement de données ou
        de statistiques
  \item Utilisation de traitements de texte (Word) pour la production de
        modèles ou de documents
  \item Gestion d'agendas, de plannings ou de réservations
  \item Accueil et réponse aux demandes récurrentes des personnels et
        étudiants
  \item Classement, archivage et gestion documentaire
  \item Opérations liées aux achats, remboursements ou finances
  \item Autre : \underline{\hspace{4cm}}
\end{itemize}

\textbf{F2.} Dans votre travail quotidien, quelles tâches vous semblent
les plus chronophages ou les plus fastidieuses ?
\textit{(Plusieurs réponses possibles)}
\begin{itemize}[label=$\square$, leftmargin=1.5cm]
  \item Rédiger ou mettre en forme des documents (courriels, rapports,
        notes de service\ldots)
  \item Rechercher une information dans plusieurs systèmes ou fichiers
  \item Répondre à des questions récurrentes (demandes fréquentes des
        personnels ou étudiants\ldots)
  \item Saisir ou organiser manuellement des données
  \item Traduire ou comprendre des contenus en langue étrangère
  \item Coordonner des réunions ou des plannings entre plusieurs
        interlocuteurs
  \item Autre : \underline{\hspace{4cm}}
\end{itemize}
 
\textbf{F3.} Effectuez-vous des opérations répétitives à intervalles
réguliers dans votre travail ? \textit{(Par exemple : générer chaque mois
le même type de tableau, envoyer chaque semestre les mêmes types de
communications\ldots)}
\begin{itemize}[label=$\square$, leftmargin=1.5cm]
  \item Oui, en grande quantité
  \item Oui, dans une certaine mesure
  \item Rarement
  \item Pratiquement jamais
\end{itemize}
 
\textbf{F4.} Avez-vous déjà entendu parler d'outils capables de rédiger
des documents, de résumer des informations ou de répondre à des questions
(tels que ChatGPT, Claude, Copilot\ldots) ?
\begin{itemize}[label=$\square$, leftmargin=1.5cm]
  \item Oui, je les utilise déjà
  \item Oui, mais je ne les ai jamais utilisés
  \item J'en ai entendu parler, mais je ne sais pas ce qu'ils peuvent
        faire concrètement
  \item Non, je ne connais pas du tout ce type d'outil
\end{itemize}
 
\textbf{F5.} Au regard des tâches que vous venez de décrire, dans quels
domaines pensez-vous qu'un outil d'\gls{ia} pourrait vous aider ?
\textit{(Plusieurs réponses possibles)}
\begin{itemize}[label=$\square$, leftmargin=1.5cm]
  \item Rédiger ou reformuler des courriels, notes et rapports
  \item Rechercher rapidement une information dans un volume important de
        documents
  \item Répondre automatiquement aux questions fréquentes
  \item Organiser ou analyser des données
  \item Traduire des documents
  \item Je ne vois pas en quoi l'\gls{ia} pourrait m'aider
  \item Je ne suis pas en mesure de répondre à cette question
\end{itemize}
 
\textbf{F6.} Votre responsable hiérarchique vous a-t-il déjà évoqué
l'utilisation d'outils d'\gls{ia} dans votre travail ?
\begin{itemize}[label=$\square$, leftmargin=1.5cm]
  \item Oui, il/elle encourage activement leur utilisation et fournit un
        accompagnement
  \item Oui, il/elle en a parlé, mais sans donner de suite concrète
  \item Non, ce sujet n'a jamais été abordé
  \item Non, il/elle y est explicitement opposé·e
\end{itemize}
 
\textbf{F7.} Avez-vous accès à des outils d'\gls{ia} générative fournis
par l'établissement ?
\begin{itemize}[label=$\square$, leftmargin=1.5cm]
  \item Oui
  \item Non
  \item Je ne sais pas
\end{itemize}
 
\textbf{F8.} Êtes-vous inquiet·ète que le développement de l'\gls{ia}
affecte votre poste ou le contenu de votre travail ?
\begin{itemize}[label=$\square$, leftmargin=1.5cm]
  \item Oui, très inquiet·ète
  \item Oui, plutôt inquiet·ète
  \item Non, plutôt pas inquiet·ète
  \item Non, pas du tout inquiet·ète
  \item Au contraire, je pense que l'\gls{ia} pourrait valoriser mon
        travail
\end{itemize}
 
\textbf{F9.} Pensez-vous que l'\gls{ia} générative pourrait améliorer
votre qualité de vie au travail (gain de temps, réduction des tâches
répétitives) ?
\begin{itemize}[label=$\square$, leftmargin=1.5cm]
  \item Oui, certainement
  \item Oui, probablement
  \item Non, probablement pas
  \item Non, certainement pas
\end{itemize}

\medskip

\subsection{Section G — Clôture}

\textbf{G1.} Quelles actions prioritaires Polytech devrait-elle mener concernant l'\gls{ia} générative ? \textit{(Plusieurs réponses possibles)}
\begin{itemize}[label=$\square$, leftmargin=1.5cm]
  \item Définir des règles claires d'usage
  \item Proposer des formations pour tous les personnels
  \item Fournir des outils institutionnels sécurisés
  \item Créer un réseau de référents \gls{ia} par département
  \item Adapter les modalités d'évaluation des étudiants
  \item Sensibiliser aux enjeux éthiques et environnementaux
  \item Autre : \underline{\hspace{4cm}}
\end{itemize}

\textbf{G2.} Avez-vous des suggestions ou des commentaires complémentaires sur l'usage de l'\gls{ia} générative à Polytech ?

\vspace{3cm}

\textbf{G3.} Accepteriez-vous d'être contacté·e pour un entretien complémentaire ou pour participer à un groupe de travail sur l'\gls{ia} ?
\begin{itemize}[label=$\square$, leftmargin=1.5cm]
  \item Oui $\rightarrow$ Votre adresse courriel : \underline{\hspace{5cm}}
  \item Non
\end{itemize}

\bigskip

\noindent\textit{Merci pour votre participation. Vos réponses contribueront à adapter l'accompagnement institutionnel aux besoins réels de la communauté Polytech.}
